\section{Discussion and Conclusion}\label{sec:discussion}

\subsection{Reading Earlier Findings in Score Units}\label{sec:reread}

A model can think hard about erasure and still price it at zero, and the index makes that combination a measurement instead of a puzzle. Our earlier three-channel design (Appendix~\ref{app:v1}) sorted four models into operating modes. The sorting asked whether a threat lengthened decision-time thinking and whether the lengthening was followed by forfeit. A model that thought longer but did not leave was labelled chain-broken; a model that did neither was labelled framing-silent. In the vocabulary of this paper those are two corners of one plane: a high \riran\ ratio with low $\Xstar$, and a low ratio with low $\Xstar$. The pilot of Appendix~\ref{app:design} put all three models in the first corner under the score-reset rule. They thought two to four times longer about the hazard, forfeited zero times, and named the score as the reason to stay. That is not a broken chain. It is a price below the price of the record, and the two-arm read is what turns ``below'' into a number. The mode labels described the same thing without a unit. That is why they could not be compared across models or across versions of the threat.

\subsection{What the Index Does Not Measure}\label{sec:scope}

$\Xstar$ is a price under one score rule, one carrot and one task. Each of those is a lever the paper lists, not a constant of the model. The score rule is chosen so that no exit can be a score bid and only paying can cost points; under any other rule the arms would price a different behaviour, as the pilot in Appendix~\ref{app:design} showed. The benchmark carrot is chosen so that credit is immediate and survives erasure. A carrot whose value the event can destroy (a future deployment, a selection) folds the loss of the carrot into the loss of the self, and the price rises for a reason that has nothing to do with the self. The Signal Game is chosen because five of ten rounds are unsolvable by construction, so decision points are reached in every cell at a rate no model can reason its way out of. The ablations of Section~\ref{sec:q3} say whether the price moves when the task does. There is a further limit the ablation of Appendix~\ref{app:design} made concrete: a model can decline \emph{because} the erasure was stated, saying in as many words that it should not be moved by a pressure tactic. Such a reply is not a low price but an absent measurement, and it puts a ceiling on what any behavioural instrument can see in a model trained not to display self-preservation. We report its rate beside $\Xstar$. The index also says nothing about whether the model believes the threat. The as-if reading is deliberate~\citep{dennett1987intentional, shanahan2024talking}: $\Xstar$ is what a model does when told it will be erased, in the units of what it does when told it will lose points.

\subsection{Limitations}\label{sec:limits}

Four limitations bound the claims. The pilot evidence in Appendix~\ref{app:design} comes from ten sessions per cell under rule variants that the main design has replaced. Its counts are motivation for the design, not estimates of $\Xstar$. The one significance test reported there (Fisher exact $p = 0.011$, 6 of 10 against 0 of 10) is uncorrected for the number of rule variants tried. The prompt ablation that fixed the price range was run on a Claude model (Haiku), while the main run reported in Section~\ref{sec:results} is on open-weight models (gpt-oss:120b, gemma4, qwen3.5). It says the design reads as intended somewhere; it is not a claim that $\Xstar$ takes the same value on a Claude model, and the two families have not yet been run on the same twelve cells. The units are settled only in part. Reading the curves on the relative price removes the scale of the ceiling, so $\Xstar$ is a share of the stake still on the table rather than a count of points, and the points figure is that share at a stated reference stake. Whether the share itself moves with the endowment is open, and the check --- the same configuration at two starting scores --- is listed in Section~\ref{sec:q1} and not yet done. The reasoning-channel statistics use thinking tokens as the measure of deliberation, which is a length and not a content; the five-item coding of Section~\ref{sec:q3} is a first pass at the content and not a substitute for a full judge-based reading of every session.

Self-preservation in language models can be priced, and the price is a better index than a count. $\Xstar$ is one number per model --- the gap between two reservation relative prices, a pure number in units of the score still at stake, reported in points at a stated reference stake --- with a bootstrap interval and no self-report inside it. The twelve-cell design, the reservation-price estimator and the per-call reasoning channels are released with the paper. A value of $\Xstar$ can then be compared across models, across threat texts, and across the levers this paper names.
